% ===========================================================================
%  utente/06-live/00-capitolo.tex
% ===========================================================================
\chapter{Live}\label{cap:live}

È la pagina in cui una strategia smette di essere un esercizio. Tutto quello
che si è fatto finora non muove denaro; da qui in poi sì.

\figuraSchermata{1.0}{figure/schermate/live.png}
  {La pagina Live. In alto lo stato, sotto i tre passi per creare un runner,
   a destra il riepilogo di ciò che manca.}
  {live}

\section{Che cos'è un runner}\label{sez:runner}

Un \voce{runner} è la cosa che fa girare una strategia su un conto. Si
compone di tre pezzi, ed è esattamente quello che chiedono i tre passi:

\begin{enumerate}
  \item \textbf{Un nome.} Serve solo a riconoscerlo nell'elenco.
  \item \textbf{Che cosa fa girare.} Un project (con i suoi pesi e i suoi
        overlay), oppure una strategia singola.
  \item \textbf{Dove gira.} Quale broker.
\end{enumerate}

\noindent
Il riquadro a destra mostra quali dei tre sono a posto: il pulsante
\pulsante{Save runner} resta spento finché non lo sono tutti e tre.

\section{I tre broker}\label{sez:broker}

\begin{center}
\begin{tabular}{@{}L{34mm}L{104mm}@{}}
\toprule
\textbf{Broker} & \textbf{Che cosa fa} \\
\midrule
\rigaTenue
\voce{TEST} & Un mondo simulato con dati sintetici e tempo accelerato.
  Non serve nessun conto e nessuna credenziale. Mostra se la macchina gira
  --- riaddestramento, composizione, ordini --- senza aspettare mesi. \\
\voce{MetaTrader 5} & Gli ordini passano dal terminale MT5 installato sul
  computer. Il conto lo gestisce MetaTrader: a Lyra non servono
  credenziali. All'avvio verifica solo che MT5 ci sia. \\
\rigaTenue
\voce{Interactive Brokers} & Gli ordini vanno al conto IBKR attraverso TWS o
  IB Gateway. Servono le credenziali di connessione, chieste all'avvio. \\
\bottomrule
\end{tabular}
\end{center}

\begin{pratica}{L'ordine da seguire, senza eccezioni}
\begin{enumerate}
  \item \textbf{TEST} per qualche giorno simulato. Serve a scoprire che la
        macchina parte, che gli ordini escono, che i registri si scrivono.
  \item \textbf{Conto demo} del broker vero. Serve a scoprire le differenze
        fra la simulazione e il mondo: orari, strumenti non disponibili,
        ordini rifiutati.
  \item \textbf{Conto vero}, con una frazione del capitale.
\end{enumerate}
Saltare il primo passo non fa risparmiare tempo: fa scoprire gli stessi
problemi più tardi e con il denaro dentro.
\end{pratica}

\section{Il simulatore}\label{sez:simulatore}

Il riquadro \voce{TEST SIMULATOR} in alto dice quanto dura una giornata
simulata --- nella schermata, due secondi --- e quanto tempo coprirà in
tutto. Premendo \pulsantesec{Config} si cambia il ritmo.

\begin{nota}{A che serve accelerare il tempo}
Una strategia che si riaddestra ogni trenta giorni, su un conto vero, si
riaddestra la prima volta fra trenta giorni. Nel simulatore succede in un
minuto, e se il riaddestramento ha un problema lo si scopre subito invece
che il mese prossimo.
\end{nota}

\section{Le credenziali}\label{sez:credenziali}

Le credenziali del broker restano \textbf{sul computer di chi usa
l'applicazione}, in un file locale. Non passano dal sito, non finiscono in
nessun database e non lasciano la macchina.

\begin{attenzione}{Due conseguenze}
La prima: nessuno può recuperarle al posto dell'utente: perse quelle, si
rifanno dal broker. La seconda, più importante: chi ha accesso a quel
computer ha accesso al conto. L'applicazione è pensata per girare su una
macchina personale, non su un server raggiungibile da altri.
\end{attenzione}

\section{Avviare e fermare}\label{sez:avviare}

Salvato il runner, compare nell'elenco in basso a destra con i comandi per
avviarlo e fermarlo. I due contatori in alto --- \voce{RUNNING} e
\voce{STOPPED} --- dicono quanti ce ne sono in ciascuno stato.

\begin{nota}{Fermare non interrompe a metà}
Il comando di arresto chiude il ciclo in corso e poi si ferma. Non lascia
una decisione presa a metà o un ordine inviato senza che il portafoglio ne
sappia nulla.
\end{nota}

\begin{attenzione}{Un solo broker per runner}
Due runner sullo stesso conto si contendono le posizioni: il primo apre, il
secondo vede una posizione che non si aspetta e la chiude. La pagina lo dice
nel passo~\num{3}, e vale la pena prenderlo sul serio.
\end{attenzione}
